\documentclass[11pt]{article}

\usepackage[margin=1in]{geometry}
\usepackage{amsmath,amssymb,amsfonts}
\usepackage[numbers,sort&compress]{natbib}
\usepackage{graphicx}

\newcommand{\Vori}{\widetilde V}
\newcommand{\Eori}{\widetilde E}
\newcommand{\start}{\mathrm{start}}
\newcommand{\Vend}{\mathrm{end}}
\newcommand{\ACO}{\mathrm{ACO}}
\newcommand{\argmin}{\mathop{\mathrm{arg\,min}}}
\newcommand{\argmax}{\mathop{\mathrm{arg\,max}}}

\title{Feasible-Walk Ant Colony Optimization for Pangenome-Guided Assembly}
\author{Anonymous Authors}
\date{\today}

\begin{document}
\maketitle

\begin{abstract}
Pangenome-guided sequence assembly reconstructs target genomes by traversing population graphs, reducing reference bias while preserving population variation. The central inference problem is Oriented Tangle Resolution: the solver must identify a valid oriented walk through the graph that agrees with copy-number estimates derived from reads. Existing Quadratic Unconstrained Binary Optimization (QUBO) formulations express this decision as time-indexed binary assignments and add penalties for invalid transitions, so graph validity is handled as an objective term instead of a property of the sample space.

We reformulate this resolution problem so that graph validity is an inherent property of every sampled candidate. Specifically, we instantiate an Ant Colony Optimization (ACO) path finder that constructs legal oriented walks for minigraph assembly and scores them with the same objective used by the original pipeline. In this path finder, the residual heuristic, termination rule, and pheromone feedback operate directly on legal choices of the next oriented state. We further train a neural transition prior on search trajectories, which biases those choices without changing the feasible action set. On held-out Graphical Fragment Assembly (GFA) files produced by the minigraph pipeline, learned logits outperform shuffled and random score controls, while a zero-prior control remains competitive on the short optimization proxy. In end-to-end minigraph assembly, the learned prior improves coverage, used sequence, and break count over the non-neural feasible-walk sampler on paired held-out assemblies.
\end{abstract}

\section{Introduction}

Genome assembly remains difficult in repetitive and structurally variable regions. De novo assemblers avoid reference bias but can fragment when reads do not bridge repeats, whereas workflows based on a linear reference can force the reconstruction toward alleles already present in that reference. Pangenome graphs with sequence orientation offer a richer substrate because they represent sequence fragments, observed adjacencies, and multiple allelic or structural alternatives in a single graph \citep{ballouz2019reference,schneider2017grch38,matthews2024gentle,liao2023draft,hickey2020vg,li2020minigraph}. Once reads from a new sample are mapped to this graph, assembly becomes a constrained path-finding problem in which the solver must choose a valid graph walk that explains read-derived evidence and can be decoded into sequence.

This inference bottleneck is visible in the pangenome-guided assembly pipeline proposed by \citet{cudby2026pangenome}. In this workflow, short reads are aligned to a population graph with annotators such as minigraph, GraphAligner, vg Giraffe, or kmer2node, and the resulting depth profiles are normalized into node-copy estimates \citep{li2020minigraph,rautiainen2020graphaligner,siren2021giraffe,cudby2026pangenome}. The optimizer must then find an oriented walk that accurately reflects these estimates. The Quadratic Unconstrained Binary Optimization (QUBO) interface encodes this decision with time-indexed binary variables and adds penalties for illegal transitions \citep{cudby2026pangenome,dunning2018mqlib}; as a result, graph validity competes with copy-number agreement inside the objective, and search effort can move through binary assignments before they decode into valid biological walks.

Because the minigraph pangenome is a compressed representation of related genomes, its construction keeps the graph broadly co-linear with the input genomes while collapsing many small local variants into representative or consensus sequences at nodes. This compression preserves the larger structural ambiguity relevant to assembly: repeats, changes in copy number, inversions, and translocations remain as graph structures that the path finder must resolve. Because the held-out target is observed through simulated short reads, the inferred walk must both explain the copy-number surface and decode into a sequence that passes consensus-quality evaluation against the target. Structural feasibility should therefore be treated as a hard constraint rather than a soft penalty.

Because short-read evidence gives noisy estimates of copy number, the optimization view remains useful. Pathfinder, the Oatk submodule used as a comparator for walk search in the original benchmark, assigns weights to the graph and searches for walks that explain them \citep{zhou2025oatk,cudby2026pangenome}. This prescriptive traversal is effective when coverage is accurate, but it can split an assembly when coverage is imperfect or when unused pangenome nodes remain in the graph. A cost-based objective can instead prefer the walk that best explains the observed copy-number surface, which means the bottleneck is not optimization itself but the search space used to represent the assembly decision.

We address this bottleneck by changing the inference formulation while preserving the biological objective and the assembly evaluation. Instead of sampling over assignments that must be penalized into validity, we sample valid oriented walks so that graph validity is guaranteed by construction and the copy-number objective can measure evidence for the assembled path directly. In our Ant Colony Optimization (ACO) instantiation \citep{dorigo1997acs,dorigo2004aco}, each ant starts from an admissible oriented state, repeatedly chooses a legal successor, and terminates at either the artificial end state or the fixed horizon inherited from the QUBO formulation. This makes transition choice the central object of search: the residual-copy heuristic scores legal successors by remaining copy-number demand, and the pheromone update reinforces complete legal walks after path-level evaluation.

We add data-driven guidance to this transition distribution through a graph-conditioned prior over the same legal successors used by the non-neural sampler. Neural guidance therefore changes transition preferences while preserving the feasible action space. Our experiments separate the effect of the feasible-walk formulation from the effect of the prior: at the optimization stage, the guided variant reduces energy on held-out GFAs from minigraph, and ablation controls show that learned logits are more useful than shuffled or random edge scores even though the zero-prior control remains a strong short-budget proxy baseline. We then run full assembly to measure how optimization improvements interact with consensus reconstruction.


\section{Related Work}
\label{sec:related-work}

\subsection{Classical graph-based sequence assembly}

Sequence assembly has long been formulated as inference over a graph whose structure is induced by read overlaps, $k$-mers, or repeat relationships. Broad surveys distinguish overlap--layout--consensus, string-graph, and de Bruijn graph assemblers, while emphasizing that repeats, uneven coverage, sequencing error, and heterozygosity turn assembly into a graph-disambiguation problem rather than a simple concatenation task \citep{miller2010assembly,nagarajan2013sequence}. Within the overlap-based tradition, string graphs represent irreducible read-overlap relationships after transitive reduction \citep{myers2005stringgraph}. By contrast, the de Bruijn tradition transforms assembly from an ordering problem over reads into an Eulerian traversal problem over $k$-mer-derived edges \citep{pevzner2001eulerian,compeau2011debruijn}; practical short-read assemblers such as Velvet then combine graph compaction, coverage filters, and paired-read information to simplify the graph \citep{zerbino2008velvet}.

Long-read assembly has shifted the balance from short-read repeat ambiguity toward noisy-overlap and haplotype-aware graph resolution. For example, Flye constructs repeat graphs from long, error-prone reads and resolves repeats using graph structure and read support \citep{kolmogorov2019flye}. Hifiasm instead uses phased assembly graphs to preserve haplotype information in high-fidelity long-read assembly \citep{cheng2021hifiasm}, while Verkko combines accurate long reads, ultra-long reads, and haplotype-specific markers to simplify a multiplex de Bruijn graph toward telomere-to-telomere diploid assemblies \citep{rautiainen2023verkko}. Together, these systems show that high-quality assembly depends critically on path extraction from a graph. Our setting uses a different graph source: a population pangenome graph annotated by target-read evidence, with inference defined as finding an oriented walk whose node visits explain copy-number estimates.

Coverage-aware path extraction also appears in strain-resolved and metagenomic assembly. Haploflow, for example, repeatedly finds coverage-supported paths through a de Bruijn unitig graph and reduces flow along selected paths to recover viral strains from mixed samples \citep{fritz2021haploflow}. This family of methods demonstrates that path and flow structure can be an effective assembly primitive. In our problem, however, the graph is an external pangenome reference annotated by target-read depth, the output is a single oriented walk through sequence-bearing pangenome states, and the objective penalizes residual node-copy mismatch rather than decomposing an assembly graph into multiple strain flows.

\subsection{Pangenome graphs and pangenome-guided assembly}

Pangenome graphs replace a single linear reference by a graph that represents sequence shared across, and variable among, a population. This representation is motivated by the reference-bias limitations of linear genomes \citep{ballouz2019reference} and has been surveyed as a general substrate for read alignment, variant calling, genotyping, visualization, and functional genomics \citep{eizenga2020pangenome,matthews2024gentle}. At human scale, graph pangenomes are now practical reference structures, as illustrated by the Human Pangenome Reference Consortium draft \citep{liao2023draft}. Toolchains such as vg, minigraph, and the Pangenome Graph Builder (PGGB) construct or operate on different pangenome graph models, with different tradeoffs between large structural variation, base-level alignment, and scalability \citep{hickey2020vg,li2020minigraph,garrison2024pggb}. In these graph models, a genome, haplotype, contig, or read alignment is naturally represented as a valid oriented path or walk through a bidirected sequence graph.

Several pangenome methods therefore perform path inference, but with different input-output interfaces. PanGenie uses a hidden Markov model over known haplotype paths and read k-mer counts to infer genotypes across variant bubbles \citep{ebler2022pangenie}. Recent integer-programming approaches similarly infer pangenome paths for genotyping by optimizing over an expanded graph and read-derived string evidence \citep{chandra2024phi}. These methods reconstruct or genotype haplotypes in a haplotype-resolved reference setting, usually by selecting among known haplotype mosaics or by enforcing a path through an acyclic or expanded graph. Pasa uses population pangenome graphs to scaffold prokaryote assemblies \citep{do2024pasa}, but it orders contigs using pangenome-supported matching scores rather than optimizing a read-depth copy-number walk through a minigraph sequence graph. The closest direct predecessor is the pangenome-guided sequence assembly framework of \citet{cudby2026pangenome}, which annotates a pangenome graph with read-derived copy-number estimates and formulates assembly as graph traversal optimization through a QUBO interface.

Building on this pangenome-guided assembly setting, we keep the biological objective and evaluation boundary fixed while changing the optimization space and search mechanism for the resolution task. Since valid graph walks are already a basic object in sequence and variation graphs, we use them as the constructive search space for path finding under copy demand. As a result, every candidate is a legal oriented pangenome walk: the heuristic measures residual demand of biological nodes, and pheromone feedback adapts a stochastic transition distribution over legal successors. This interface differs from hidden Markov model (HMM) or integer-programming haplotype inference, and it also differs from scaffolding-by-matching and from QUBO search over time-indexed binary assignments that must be penalized into validity.

\subsection{Heuristics and meta-heuristics for sequence assembly}

Before modern graph assemblers became dominant, the DNA fragment assembly problem was frequently studied as a combinatorial optimization problem over fragment permutations, overlap scores, and contig structure. Genetic algorithms were explored early for the layout phase of fragment assembly \citep{parsons1995genetic}, and subsequent work compared genetic algorithms, scatter search, local-search variants, and hybrid evolutionary search for the fragment assembly problem \citep{luque2005metaheuristics,alba2008hybrid}. Ant-colony approaches were also applied to shotgun fragment assembly: Meksangsouy and Chaiyaratana formulated fragment ordering using ant-colony-system search \citep{meksangsouy2003dna}, while Wetcharaporn et al. combined Smith--Waterman overlap computation, ant-colony fragment ordering, and nearest-neighbor contig assembly \citep{wetcharaporn2006dna}. Later benchmark and optimization studies standardized comparisons for DNA fragment assembly and explored hybrid metaheuristic overlap--layout--consensus pipelines \citep{mallen2013benchmark,uzma2020optimizing}.

These works are important precedents because they treat assembly as a search problem rather than only as a deterministic graph-cleaning pipeline. Their search spaces mainly optimize read or fragment orderings under overlap-based objectives. Ant Colony Optimization itself is a general metaheuristic for discrete graph search in which stochastic constructive agents are biased by pheromone memory and local heuristic information \citep{dorigo1997acs,dorigo2004aco,blum2005aco}. The assembly space studied here is different: ants traverse an oriented pangenome graph, illegal transitions are absent from the action set, and the terminal objective is copy-number agreement with target-read evidence.

\subsection{Learned and neural sequence assembly}

Learning-based assembly methods have recently begun to target graph resolution directly. GNNome formulates \emph{de novo} assembly as path identification in an assembly graph and uses geometric deep learning to predict graph edges or paths corresponding to reconstructed sequences \citep{vrcek2025gnnome}. This line of work was preceded by an arXiv version showing that a graph convolutional model trained on simulated CHM13-derived graphs could generalize to other human chromosomes and improve contiguity over hand-crafted heuristics on the same graphs \citep{vrcek2022untangle}. DipGNNome extends the geometric-learning view toward diploid assembly by combining haplotype-aware assembly graphs, supervised edge labels, and beam-search decoding \citep{schmitz2025dipgnnome}. Reinforcement-learning approaches have also been studied for genome assembly, but recent Q-learning analyses report poor scalability and difficulty obtaining competitive assembly quality in larger environments \citep{padovani2025rlassembly}.

The learned prior is related to neural assemblers because it learns graph-local decisions for path extraction, while using a different input distribution and inference interface. GNNome and DipGNNome operate on \emph{de novo} assembly graphs built from the target reads and learn to decode contigs from those graphs. Here, the graph is a minigraph pangenome with copy numbers derived from read depth, and the learned model enters as a transition prior inside the walk optimizer. Because the neural component ranks legal adjacencies rather than defining the graph itself, every guided rollout remains valid and is evaluated through the same objective and downstream parser.

\subsection{Learning-guided combinatorial search}

Our learned transition prior also connects to neural combinatorial optimization and learning-guided search. Pointer networks and reinforcement-learning formulations established that neural policies can construct solutions to routing-like combinatorial problems \citep{vinyals2015pointer,bello2017nco}. Attention models, multiple-start policy optimization, and simulation-guided beam search further improved neural construction and decoding for routing benchmarks \citep{kool2019attention,kwon2020pomo,choo2022simulation}. More generally, graph neural networks have been used to learn heuristics over graph-structured optimization states \citep{khalil2017learning,kipf2017gcn,hamilton2017graphsage,velickovic2018gat}. Search-guided learning has also been successful in domains where a learned policy or value function biases a symbolic search process rather than replacing it outright \citep{kocsis2006uct,coulom2006mcts,silver2016alphago,silver2018alphazero,orseau2021phs}.

Our learned ACO variant follows the same hybrid principle: a neural model shapes a classical search distribution, while the classical procedure maintains feasibility, stochastic exploration, and adaptive memory. The design is closest in spirit to dynamic neural guidance for ACO \citep{tran2026dynaco}, but the biological problem imposes constraints absent from routing benchmarks. A transition is useful only if it is a legal oriented adjacency in the pangenome, and repeated traversal of its biological node should also improve agreement with estimated copy. We therefore learn an edge prior aligned with the walk structure of pangenome-guided assembly.

\section{Pangenome-Guided Assembly as Feasible-Walk Inference}

\subsection{Assembly pipeline and inference object}

The assembly instance begins with two inputs: a Graphical Fragment Assembly (GFA) pangenome and short reads from a target sample \citep{li2016miniasm,li2020minigraph}. In the GFA, segment records store DNA strings, whereas link records describe adjacencies observed in the population graph. Because the target sample supplies reads rather than a known path, the synthetic benchmark first generates related haploid genomes, uses one subset to build the pangenome, and shreds held-out genomes into reads so that the final assembly can be compared with known truth.

Read-to-graph annotation turns these raw inputs into the weighted graph seen by the path-finding algorithm. The original benchmark studies minigraph, GraphAligner, vg Giraffe, and kmer2node as annotation sources, which differ in whether they supply evidence from read depth, path alignment, or k-mer support \citep{li2020minigraph,rautiainen2020graphaligner,siren2021giraffe,cudby2026pangenome}. The experiments here use the slice produced by the minigraph pipeline, where \texttt{SC:f} depth tags are normalized into copy-number estimates before optimization. The resulting input is an annotated pangenome with sequence-bearing nodes, graph adjacencies, node lengths, and a nonnegative copy-number surface estimating how many times each biological segment should be traversed by the target genome.

In our experiments, the intervention is the path-finding stage of the standard assembly pipeline, whose other stages are graph construction, read simulation, mapping, normalization, and evaluation. Our solver replaces path finding while upstream annotations and downstream consensus metrics remain unchanged, allowing us to isolate the inference formulation from changes to graph construction or sequence evaluation.

\begin{figure}[t]
    \centering
    \includegraphics[width=\linewidth]{images/aco_pipeline.jpg}
    \caption{Pangenome-guided assembly pipeline with the feasible-walk path-finding module. Graph construction, read simulation, mapping, normalization, and evaluation remain fixed while the path-finding component changes.}
    \label{fig:aco_pipeline}
\end{figure}

The path-finding stage receives a graph that may contain unused population nodes, repeated sequence, collapsed local variation, and noisy node depths. If annotated copy numbers are treated as hard traversal instructions, a zero- or low-coverage node can split the output even when a nearby path through that node would produce a better candidate sequence. Treating copy-number agreement as a cost gives the solver room to prefer a single coherent walk through a noisy region when the estimated copy-number surface is imperfect.

We evaluate the output of path finding both as an optimization solution and as an assembled sequence. To obtain metrics at the sequence level, we first convert the oriented node list into a Graphical mApping Format (GAF)/path representation \citep{li2020minigraph}; we then concatenate node sequences with reverse complementation when necessary and align the candidate sequence to the known synthetic truth. These metrics capture different failure modes: coverage measures how much of the true target is recovered, used sequence measures how much of the candidate is supported by the truth, contigs and N50 describe continuity, breaks count false joins or large path errors, and indels, difference blocks, and identity measure finer sequence accuracy. Because the downstream parser is part of the evaluation pipeline, an improved path objective is meaningful only when it also improves these metrics for assembly quality.

\subsection{Annotated graph problem}

The input is an annotated GFA-like pangenome graph
\begin{equation}
G=(V,E,w,L,s),
\end{equation}
where $V$ is the set of biological sequence nodes, $E$ is the set of graph adjacencies, $w:V\rightarrow\mathbb R_{\ge 0}$ gives the normalized estimate of copy number derived from read-depth tags, $L(v)$ is node length, and $s(v)$ is the DNA sequence stored in node $v$. When evidence from reads or k-mers supplies traversal counts for edges, we denote the corresponding oriented count by $a(i,j)$. The solver output is an ordered walk through graph states, decoded into DNA by concatenating the corresponding node sequences with reverse complementation when needed.

The un-oriented form of the problem asks for a walk $W$ whose visit count $\#_W(v)$ matches the estimated copy number $w(v)$. Following the pangenome-guided assembly formulation of \citet{cudby2026pangenome}, this copy-number path problem is Tangle Resolution:
\begin{equation}
J_{\mathrm{TR}}(W)=\sum_{v\in V}\left(\#_W(v)-w(v)\right)^2.
\end{equation}
This objective captures the cost-based path-finding idea, but it does not yet account for the double-stranded nature of genomic sequence.

Oriented Tangle Resolution is the corresponding problem on a double-stranded sequence graph. For every biological node $v$, the solver creates two oriented states $v_+$ and $v_-$. The oriented graph is
\begin{equation}
\widetilde G=(\Vori,\Eori),\qquad \Vori=\{v_+,v_-:v\in V\}\cup\{\Vend\}.
\end{equation}
Every GFA link induces a directed transition together with its reverse-complement mate. To allow a walk to terminate before the fixed horizon, the augmented graph includes the artificial end state $\Vend$. The horizon follows the same copy-number scaling used by the assembly benchmark:
\begin{equation}
T=\max\left\{\left\lfloor \kappa \sum_{v\in V} w(v)\right\rfloor,1\right\},
\label{eq:horizon}
\end{equation}
with the same horizon scale $\kappa$ used by the assembly benchmark.

Given a walk $W$, let $c_W(v)$ be the number of visits to biological node $v$. In the oriented representation,
\begin{equation}
c_W(v)=\#_W(v_+)+\#_W(v_-).
\label{eq:oriented-count}
\end{equation}
The Oriented Tangle Resolution objective is the node-copy loss
\begin{equation}
J_V(W)=\sum_{v\in V}\left(c_W(v)-w(v)\right)^2.
\label{eq:node-objective}
\end{equation}
The solver uses Eq.~\eqref{eq:node-objective} as its path-finding target. We report the oriented setting because it preserves the double-stranded biological structure while avoiding the allele-assignment ambiguity of Diploid Tangle Resolution. Although diploid and polyploid extensions would replace the single walk by multiple walks whose combined oriented visit counts explain the same surface, the experiments here use the single-walk haploid setting evaluated in the minigraph assembly comparison.

\subsection{QUBO interface for Oriented Tangle Resolution}

After decoding, candidate walks are compared under the same objective and the same downstream parser used to score paths extracted from the QUBO interface. This compatibility keeps GAF conversion, sequence extraction, and evaluation of consensus quality fixed while removing repair of invalid transitions from the solver.

The QUBO interface takes a different route by converting the problem into a time-indexed binary assignment. For each time $t\in\{1,\ldots,T\}$ and each oriented state $u\in\Vori$, the encoding introduces a binary variable
\begin{equation}
x_{t,u}=
\begin{cases}
1, & \text{state }u\text{ is selected at time }t,\\
0, & \text{otherwise}.
\end{cases}
\end{equation}
A valid walk requires exactly one active state at each time and a legal graph transition between consecutive active states. Let $\mathcal E_{\mathrm{end}}$ be the legal transition set after adding termination edges into $\Vend$ and the self-loop $(\Vend,\Vend)$, which lets a completed walk remain at the end state. The QUBO objective combines three terms:
\begin{align}
C_1(x)
&=
\Lambda_1\sum_{t=1}^{T}
\left(\sum_{u\in\Vori}x_{t,u}-1\right)^2,
\label{eq:qubo-one-state}\\
C_2(x)
&=
\Lambda_2\sum_{t=1}^{T-1}
\sum_{\substack{u,v\in\Vori\\(u,v)\notin\mathcal E_{\mathrm{end}}}}
x_{t,u}x_{t+1,v},
\label{eq:qubo-transition}\\
C_3(x)
&=
\sum_{v\in V}
\left(\sum_{t=1}^{T}(x_{t,v_+}+x_{t,v_-})-w(v)\right)^2.
\label{eq:qubo-copy}
\end{align}
Here $C_1$ enforces one selected oriented state per time step, $C_2$ penalizes illegal transitions, and $C_3$ is the oriented node-copy loss. Because the benchmark horizon scales with the total estimated copy, the walk can revisit nodes when doing so improves global agreement. The virtual end state allows early termination without forcing unused positions to carry sequence.

The QUBO solved by MQLib or other binary optimizers \citep{dunning2018mqlib} is the weighted sum
\begin{equation}
C_{\mathrm{QUBO}}(x)=C_1(x)+C_2(x)+C_3(x),
\label{eq:qubo-total}
\end{equation}
up to the standard conversion into a symmetric quadratic matrix. Decoding a feasible assignment gives the oriented state sequence $(u_1,\ldots,u_T)$, with trailing end states removed before GAF/path conversion and sequence extraction. The same decoded path-level loss in Eq.~\eqref{eq:node-objective} is then used to compare candidates produced by the QUBO interface and by our sampler.

The time-indexed representation gives the original MQLib backend \citep{dunning2018mqlib} a standard input format while changing the search problem. The oriented encoding uses $(2|V|+1)T$ binary variables for a single walk, and the diploid form introduces a second assignment over time. A binary optimizer therefore explores a large space in which penalties push assignments toward feasibility. Penalty weights must be large enough to enforce constraints, yet small enough to preserve numerical differences in copy agreement. Our formulation moves the validity constraint from the objective into the sample space: the scoring rule and downstream parser remain unchanged, and every candidate is legal at construction time.

The reformulation also exposes a direct route for edge-aware evidence, an important limitation of the original QUBO route. Annotation pipelines can produce edge support when reads or k-mers traverse from one node to another, and such evidence can help distinguish orientation choices that visit the same nodes but use different adjacencies. If $e_W(i,j)$ counts how many times $W$ traverses oriented edge $(i,j)$, an edge-copy loss can be added as
\begin{equation}
J_E(W)=\sum_{(i,j)\in \Eori}\left(e_W(i,j)-a(i,j)\right)^2,
\label{eq:edge-objective}
\end{equation}
and the terminal score becomes $J(W)=J_V(W)+\lambda_EJ_E(W)$. In a time-indexed QUBO, $e_W(i,j)$ corresponds to products $x_{t,i}x_{t+1,j}$; squaring summed edge residuals then creates higher-order interactions that call for auxiliary variables. In ACO, by contrast, the same quantity is a counter over sampled transitions. The main minigraph comparison keeps the node-copy objective to match the original MQLib comparison, while this extension identifies how the walk sampler can incorporate transition-level read evidence while preserving the candidate representation.

\paragraph{Feasibility property.}
For any ant trajectory $W=(u_0,\ldots,u_\ell)$ produced by the sampler, each nonterminal transition $(u_t,u_{t+1})$ belongs to $\Eori$ and each state belongs to $\Vori$. The trajectory is therefore valid by construction. The objective can still score poor, overlong, or copy-mismatched walks, but it no longer has to repair illegal graph transitions. Feasibility is enforced by the transition kernel; quality remains measured by the same loss and assembly evaluation.

\section{Feasible-Walk Ant Colony Optimization for Oriented Tangle Resolution}

\subsection{Ant construction}

The feasible-walk sampler constructs a batch of complete legal walks. Let $S(u)$ be the set of legal successors from oriented state $u$, including $\Vend$ when termination is allowed. At step $t$, an ant at state $u_t$ samples $u_{t+1}\in S(u_t)$ according to
\begin{equation}
p(u_{t+1}=v\mid u_t,c_t)
=
\frac{\tau_{u_tv}^{\alpha_\tau} \eta(v,c_t)^{\beta_\eta}}
{\sum_{z\in S(u_t)} \tau_{u_tz}^{\alpha_\tau}\eta(z,c_t)^{\beta_\eta}},
\label{eq:aco-transition}
\end{equation}
where $\tau_{uv}$ is pheromone, $c_t$ is the current node-visit count vector, and $\eta(v,c_t)$ is a residual-copy heuristic. The heuristic gives high score to successors whose underlying biological node still has unmet copy-number demand:
\begin{equation}
\eta(v,c_t)=
\begin{cases}
0.01+\max(0,w(b(v))-c_t(b(v)))\left(1+\frac{\sqrt{1+L(b(v))}}{10}\right),
& v\neq \Vend,\\
\eta_{\mathrm{end}}(c_t),
& v=\Vend,
\end{cases}
\label{eq:residual-heuristic}
\end{equation}
where $b(v)$ maps an oriented state to its biological node. The end-state score is
\begin{equation}
\eta_{\mathrm{end}}(c_t)=
\begin{cases}
1, & \sum_{x\in V}\max(0,w(x)-c_t(x))=0,\\
0.01, & \text{otherwise}.
\end{cases}
\end{equation}
The end state therefore receives a high score only after residual demand is exhausted; otherwise it remains available but discouraged.

Legal-successor sampling gives the assembly solver two properties at once: every ant path can be converted back into a graph path directly, and residual-copy scores change as the ant consumes node demand. The same edge can therefore be attractive early in a walk but unattractive after repeated visits to its biological node.

\subsection{Pheromone feedback}

Pheromone feedback turns batches of legal walks into an adaptive search distribution by using decoded path-level scores to identify elite traces after each ant batch. These traces then update pheromone, while evaporation is applied as
\begin{equation}
\tau_{uv}\leftarrow \max\{(1-\rho)\tau_{uv},\epsilon\},
\label{eq:evaporation}
\end{equation}
where $\rho$ is the evaporation rate and $\epsilon$ prevents zero-probability transitions. The best quarter of the ants deposit additional pheromone on their traversed edges:
\begin{equation}
\tau_{uv}\leftarrow \tau_{uv}+\Delta_k
\quad\text{for each edge }(u,v)\text{ in elite ant }k.
\end{equation}
The deposit $\Delta_k$ is scaled by the ant's energy relative to the worst energy in the batch, so lower-energy ants reinforce their edges more strongly while evaporation and residual demand preserve exploration.

\paragraph{Adaptive feasible-walk distribution.}
The sampler maintains an incumbent walk $W^\star$ and a pheromone field initialized to $\tau_e=1$ on every legal transition. During one inference update, it draws $M$ legal walks from Eq.~\eqref{eq:aco-transition}, evaluates them with the decoded path-level objective, stores any lower-energy incumbent, and shifts the transition distribution through evaporation and elite-trace deposition. Additional inference compute increases the number of feedback cycles in which sampled walk quality can reshape the transition distribution. The minimum-update requirement guarantees at least one such feedback cycle.

\subsection{Compute-quality tradeoff}

Sampled walks and pheromone updates define the compute-quality tradeoff of the feasible-walk formulation. Pangenome graph builders produce subgraphs with different sizes and branching structures, so the relevant question is how effectively additional inference compute improves already-valid oriented walks rather than how often invalid assignments can be repaired after optimization.

\section{Learned Transition Priors}

\begin{figure}[t]
    \centering
    \includegraphics[width=\linewidth]{images/pipeline_sec4.png}
    \caption{Learned transition-prior interface for feasible-walk ACO.
        \textbf{Steps 1--3:} The current ant state $c_t$,
        oriented pangenome graph $\tilde{G}$, and read-derived
        copy-number evidence $w$ are combined as input.
        \textbf{Step 4:} The graph neural network computes edge logits
        $z_\theta(e;\tilde{G},w)$ \emph{once} per outer iteration
        as a held prior.
        \textbf{Steps 5--6:} The transition score
        $\log s(e \mid c_t) = \alpha\log\tau_e +
        \beta\log\eta(v,c_t) + \gamma r_\theta(e;\tilde{G},w)$
        combines pheromone, residual-copy heuristic, and neural
        prior to rank legal outgoing edges.
        \textbf{Step 7:} Inner ACO rollouts sample feasible walks,
        evaluate node-copy energy $J(W_a)$, update pheromone
        via evaporation and elite deposit (Eq.~16), and store
        trajectory snapshots $\tau^{(h,t)}$.
        \textbf{Step 8:} Online policy-gradient training replays
        stored trajectories using the saved $\tau$ snapshot,
        computes batch-centered advantages $A_a$, minimizes the
        REINFORCE loss $\mathcal{L}(\theta)$ \citep{williams1992reinforce},
        and updates the graph neural network.}
    \label{fig:pipeline_sec4}
\end{figure}

\subsection{From static guidance to dynamic transition shaping}

We align the learned signal with the search state in which it is used. A fixed heatmap can help at initialization, but the feasible-walk distribution changes as pheromone accumulates and as the incumbent improves. We therefore adapt the DyNACO idea of dynamic neural guidance \citep{tran2026dynaco} by training on decisions sampled inside the iterative search trajectory.

In the pangenome instantiation, we preserve the sampler for legal walks, the heuristic for residual copy, and the pheromone dynamics before adding a graph neural prior on edges to the transition score. Because the environment already adapts through pheromone updates and incumbent improvements, the neural policy complements search by supplying edge-level guidance to the transition distribution over feasible walks. Thus, even as feasible walks replace traveling-salesman-problem (TSP) tours and copy-number loss replaces route length, the learned object retains the role of dynamic neural guidance as an intervention on the transition distribution.

\subsection{Transition prior trained on trajectories}

We train the learned prior at the same temporal granularity at which it will influence search. A neural decision shapes a block of sampling over feasible walks, after which pheromone feedback changes the search state before the next neural decision. At each outer update, the model observes the oriented graph instance and emits a vector of edge logits once; the inner environment then runs several ACO batches with those logits held fixed while pheromone changes after every batch. The neural action is therefore a fixed-duration macro-action that biases a block of legal successor decisions before stochastic sampling, path scoring, and elite feedback produce the next macro-state.

For training graph $G$, outer update $h$, and inner batch $s$, the rollout state can be written as
\begin{equation}
\mathcal S_{h,s}=(G,w,\tau^{(h,s)},c_t,u_t,T),
\end{equation}
where $\tau^{(h,s)}$ is the pheromone vector at the start of the batch, $u_t$ and $c_t$ are the current ant state and visit counts collapsed across both orientations during construction, and $T$ is the horizon. At the outer update, the learned action is an edge-logit vector
\begin{equation}
z_\theta(G,w)=\{r_\theta(u,v;G,w):(u,v)\in\mathcal E_{\mathrm{end}}\}.
\end{equation}
The vector of edge logits is not recomputed from the live pheromone field inside the inner rollout. Instead, training on trajectories exposes the prior to changing pheromone states across rollout blocks, while deployment holds the current edge prior fixed within each block. Adaptation inside a block comes from pheromone feedback and the heuristic for residual copy; adaptation across blocks comes from the next update of the neural prior.

Deployment combines the held prior with the current ACO state at every legal successor decision. For an ant at state $u_t$ with legal successor $v$, the deployed score is
\begin{equation}
\log s(v\mid u_t,c_t,h,s)
=
\alpha_\tau\log \tau^{(h,s)}_{u_tv}
+\beta_\eta\log \eta(v,c_t)
+\gamma r_\theta(u_t,v;G,w).
\label{eq:dynaco-score}
\end{equation}
The transition distribution is the softmax of Eq.~\eqref{eq:dynaco-score} over $v\in S(u_t)$. Setting $\gamma=0$, or using a zero prior, recovers the non-learned ACO variant; consequently, the learned component is a bias over legal transitions that preserves feasible-path construction.

At deployment, $z_\theta(G,w)$ is computed once for the outer update. Then, each of the $\mathrm{mini\_H}$ inner batches samples $n_{\mathrm{ants}}$ feasible walks, scores them with the decoded path-level objective, updates pheromone from elite traces, and retains the incumbent. Evaluation uses the same loop, except that the learned checkpoint supplies the held prior and pheromone is initialized locally for each instance.

\subsection{Graph neural edge prior}

The graph neural prior estimates which legal transitions remain useful by combining the local construction state with the global copy-number surface. Its representation must distinguish residual demand from over-use, recognize the current oriented state, and compare legal successors in the context of graph density and remaining horizon. We encode these signals with three feature groups: node features include copy number, visit count, residual demand, length, orientation, current-state status, zero-copy status, and local degree; edge features include source-target incidence, structure around the end state, artificial transitions from start or to end, self-transitions, endpoint copy numbers, and residual state at the target; and global features include prefix depth, total copy-number mass, residual/over/under-use, and graph density.

The online trainer builds the prior tensor from the graph tensor at the start state, using zero visit counts, the artificial start state, and the walk horizon derived from QUBO. From this tensor, the model emits one logit for every legal edge in the augmented transition set. During replay and sampling, $\eta(v,c_t)$ changes with the ant prefix and pheromone changes across inner batches, while the neural logit for an edge from source to target remains the held prior value produced at the outer update.

Node and edge features are linearly embedded before passing through a stack of edge-aware message-passing layers. Within each layer, gated source-node messages pass through edge states, after which node embeddings are updated by a residual normalized transformation and edge embeddings are updated from the source embedding, target embedding, and previous edge state. A scalar edge head then decodes each legal transition logit from the edge state, source embedding, target embedding, and global features. The reported configuration uses 2 message-passing layers, hidden dimension 16, AdamW optimization \citep{loshchilov2019adamw}, and gradient clipping.

We insert the learned prior into a search distribution that already contains pheromone feedback and the residual-copy heuristic, so its role is to supply graph-context preferences among legal successors. It can reinforce an edge with low pheromone when the successor helps satisfy residual demand, suppress an over-reinforced edge that repeatedly creates copy-number mismatch, or separate alternatives with similar local heuristic scores but different graph context. In the paired comparison between the base sampler and the guided variant, we isolate this term in Eq.~\eqref{eq:dynaco-score} while keeping graph instances, preprocessing, number of sampled walks, objective, and consensus evaluation fixed.

\subsection{Online rollout training}

We train the prior online from feasible-walk ACO rollouts generated directly by the search procedure. For each training graph, pheromone starts at one on every legal transition; at outer update $h$, the trainer computes a detached prior $z_{\theta_{\mathrm{old}}}(G,w)$ and uses it to sample $\mathrm{mini\_H}$ inner ant batches. Before each inner batch, the trainer stores the current pheromone vector, and the C++ sampler records each ant trace as start offsets and chosen edge ids. After batch scoring, elite traces update pheromone before the next batch is sampled.

After sampling, trajectory replay evaluates the sampled decisions under the current model. Replay recomputes the log-probability of each chosen edge using the stored pheromone snapshot, the current model's held prior, and the dynamic residual heuristic reconstructed from the prefix visit counts. For an ant walk $W_a=(u_0,\ldots,u_\ell)$, the replayed log-probability is
\begin{equation}
\log p_\theta(W_a\mid \tau,z_\theta)
=
\sum_{t=0}^{\ell-1}
\log
\frac{\exp(\log s(u_{t+1}\mid u_t,c_t,h,s))}
{\sum_{v\in S(u_t)}
\exp(\log s(v\mid u_t,c_t,h,s))}.
\label{eq:replay-logp}
\end{equation}
The trajectory score is length-normalized by the number of stochastic decisions in the trace before forming the batch loss. If ant $a$ in batch $B$ produces walk $W_a$ with energy $J(W_a)$, the batch-centered cost advantage is
\begin{equation}
A_a=\bar J_B-J(W_a),
\qquad
\bar J_B=\frac{1}{|B|}\sum_{b\in B}J(W_b).
\end{equation}
The minimized replay loss is
\begin{equation}
\mathcal L(\theta)
=
\frac{1}{|B|}\sum_{a\in B}
\left(J(W_a)-\bar J_B\right)
\log p_\theta(W_a\mid \tau,z_\theta),
\label{eq:policy-loss}
\end{equation}
with stop-gradient costs. Under gradient descent, lower-energy walks have negative centered cost and therefore receive increased probability, while higher-energy walks are suppressed. The objective is a REINFORCE trajectory replay loss with batch baselines and gradient clipping \citep{williams1992reinforce}. Pheromone is treated as rollout-local search state; the transferable learned object is the neural transition prior.

We train on trajectories by optimizing the policy on walks sampled inside the pheromone loop. Because pheromone concentration, elite traces, and residual patterns change during search, the same graph can induce different training signals at different outer updates. The prior is therefore trained to complement ACO's adaptive memory: neural logits provide transition bias from graph context, while pheromone feedback supplies adaptation local to the rollout.

We draw the training distribution from the minigraph assembly pipeline, with each training instance coming from the same upstream simulation used for evaluation. Synthetic haplotypes define a population graph, a held-out genome is shredded into simulated reads, and minigraph writes an annotated query GFA with \texttt{SC:f} read-depth tags. The data-collection run stops at this annotated-GFA stage: the simulator is invoked in data-only mode, so the required \texttt{aco} solver argument serves as a placeholder for the script interface and path-finding is deferred to the training and evaluation routines. When an accepted GFA is materialized as an optimization instance, we normalize its tags into node copy numbers using the benchmark conversion, using minimum copy 0.45, integer mode, offset 0.4, and depth divisor 5.0. The learned policy therefore sees the same kind of normalized surface used by the minigraph benchmark and learns from search trajectories rather than solver-output labels.

We train on fresh online minigraph batches and evaluate on held-out minigraph instances. In each epoch, the training stream generates 32 fresh accepted GFAs from minigraph seeds and trains once on that batch; after training, the learned-prior checkpoint is used as a fixed MG coverage prior in the reported optimization and assembly comparisons. This design exposes the policy to varied graphs and demand surfaces during learning while keeping the reported paired rows on held-out minigraph data.

Because objective evaluation materializes a dense QUBO matrix indexed by time, we use a scale guard for QUBO-stage supervision: memory demand grows with both node count and horizon. The corpus index first filters generated GFAs by segment count and a cheap segment-count horizon proxy; exact benchmark-style copy numbers are then computed when a selected GFA is built as an optimization instance. This keeps data collection bounded while avoiding unnecessary copy-number conversion over discarded generated candidates. The current minigraph-specific training guard keeps GFAs with at most 80 segments and proxy horizon at most 300. Importantly, the target assembly distribution remains the distribution produced by the minigraph pipeline; the guard is an implementation constraint for dense training instances rather than a change to the upstream generator.

\begin{table}[t]
\centering
\caption{Learned-prior protocol for trajectory-aware training and paired comparison.}
\label{tab:dynaco-settings}
\begin{tabular}{ll}
\hline
Quantity & Setting \\
\hline
Training epochs & 10 \\
Training steps per epoch & 32 fresh GFAs \\
Training source & Data-only minigraph GFA stream, fresh per epoch \\
Evaluation source & held-out minigraph GFAs disjoint from the online training batches \\
Copy-number source & Copy numbers normalized from GFA \texttt{SC:f} depth tags \\
Feasibility guard & At most 80 segments and proxy horizon at most 300 \\
Online guidance steps $H$ & 10 \\
Inner steps $\mathrm{mini\_H}$ & 10 \\
Total ant batches $H\times\mathrm{mini\_H}$ & 100 \\
Ants $n_{\mathrm{ants}}$ & 32 \\
Graph neural network depth / hidden units & 4 / 64 \\
Optimizer / learning rate & AdamW \citep{loshchilov2019adamw} / $5\times10^{-4}$ \\
Gradient clipping & 1.0 \\
Pheromone exponent $\alpha_\tau$ & 1.0 \\
Residual heuristic exponent $\beta_\eta$ & 1.0 \\
Evaporation $\rho$ & 0.1 \\
Elite fraction & 0.25 \\
Neural prior weight $\gamma$ & 1.0 \\
\hline
\end{tabular}
\end{table}

\section{Evaluation Interface}

The evaluation changes the path-finding stage of the assembly pipeline. Upstream graph generation, read simulation, read-to-graph annotation, and copy-number estimation remain unchanged within each evaluated condition; downstream, the selected walk is converted to GAF/path output, decoded into sequence, and passed through the same consensus-quality evaluation. The resulting metrics measure complementary aspects of assembly quality: coverage measures truth recovery, used sequence measures support for the candidate assembly, contigs and N50 measure continuity, breaks measure false joins, and indels, difference blocks, and identity measure sequence-level accuracy. This boundary lets us test two inference spaces for the same resolution problem: penalty-constrained binary assignments and feasible oriented walks.

We separate paired evaluation from online learning to keep training variation distinct from solver comparison. Specifically, we use a fixed minigraph corpus for solver comparisons on identical graph instances, fresh per-epoch training graphs from the same minigraph pipeline, and held-out test sets for the reported optimization and assembly comparisons. In full assembly comparisons, the data-only minigraph output for each seed is generated once and then reused across solvers, so differences come from path finding rather than from independently sampled target genomes.

For assembly comparisons, we follow the original benchmark's Pathfinder preprocessing and construction of the QUBO subgraph whenever that preprocessing mode is part of the evaluated solver condition for minigraph \citep{cudby2026pangenome}. In that end-to-end condition with a Pathfinder graph, the original assembly helper converts subgraph depth to copy number. For training and standalone QUBO-stage comparisons of the learned prior, we stop the same minigraph pipeline after annotated query GFAs are produced, normalize their \texttt{SC:f} depth tags into copy numbers in the benchmark style, and expose the resulting instances for training or held-out optimization. To include the biological walk baseline in the same objective table, we run Pathfinder on these 10 GFAs and rescore its emitted paths with Eq.~\ref{eq:node-objective}. The original benchmark also studies GraphAligner, vg Giraffe, and kmer2node annotation routes; the learned-prior experiments reported here use the minigraph slice so the model sees one consistent distribution of annotated GFAs. Graphs from this pipeline contain read-depth tags, mapping noise, and preprocessing-induced subgraphs, which makes them more representative of the minigraph assembly benchmark than simple line-and-shortcut generators. Consequently, the learned prior is trained on the same kind of copy-number surface used by the minigraph assembly benchmark, while the end-to-end comparison preserves the original assembly route.

\section{Experiments}

\subsection{Evaluation protocol}

We organize the evaluation around the two contributions. First, we isolate legal-walk sampling as the non-learned carrier for the assembly formulation, then use held-out optimization instances to measure the effect of edge-level guidance against the same ACO distribution over feasible walks and against non-neural baselines. In an ablation, we vary the prior source and neural weight to test whether the neural signal carries edge-specific information beyond zero, shuffled, or random logits. Finally, we pass both walk-based solvers through the original sequence parser to measure how intermediate optimization quality becomes recovery and structural correctness at the assembly level.

Across comparisons, we preserve the biological objective and the evaluation boundary: upstream graph construction, read mapping with minigraph, normalization, path extraction, and consensus-quality evaluation are shared, so differences come from the inference formulation and its neural guidance. At the optimization stage, the comparison has three roles. First, MQLib serves as the direct external baseline because it is the QUBO heuristic backend used for Oriented Tangle Resolution in the original pangenome-guided assembly interface \citep{cudby2026pangenome,dunning2018mqlib}. Second, Pathfinder represents the original walk-search baseline in that pipeline, combining node and edge weights in an exhaustive traversal designed for coverage-annotated assembly graphs \citep{zhou2025oatk,cudby2026pangenome}. Finally, beam search and A* search provide graph-search controls adapted to the same oriented objective \citep{lowerre1976harpy,hart1968astar}, while feasible-walk ACO provides the non-neural internal variant that isolates the learned guidance layer \citep{dorigo1997acs,dorigo2004aco}.

We choose the end-to-end baseline set to match the task definition of pangenome-guided sequence assembly, where the solver receives an annotated minigraph instance and returns a copy-number-scored oriented walk for the same consensus parser. Related pangenome tools use different input-output interfaces: Pasa uses population pangenome graphs for scaffolding prokaryote assemblies \citep{do2024pasa}; PanGenie and integer-programming pangenome inference methods perform genotyping or haplotype inference from known graph structure and read evidence \citep{ebler2022pangenie,chandra2024phi}; and neural graph decoders such as GNNome and DipGNNome operate on \emph{de novo} assembly graphs \citep{vrcek2025gnnome,schmitz2025dipgnnome}. We therefore use the original QUBO/MQLib route under the same pangenome-guided assembly pipeline as the direct empirical baseline. At full assembly scope, we evaluate feasible-walk ACO and learned-prior ACO through the consensus parser on minigraph instances to measure downstream transfer from path optimization to sequence recovery.

\subsection{Ablation: feasible-walk path finding}

In the reformulation check, feasible-walk ACO and MQLib share minigraph graph generation, normalization, path extraction, and the consensus parser; the controlled change is where feasibility is enforced, either through QUBO penalties or through legal-walk construction.

To keep this comparison local, we hold assembly instance construction fixed while changing the inference space. We use three seeds with one test sequence per seed, Pathfinder graph preprocessing, and a five-second path-finding call in which the feasible-walk sampler draws 8 ants and performs at least one pheromone-update iteration.

\begin{table}[t]
\centering
\small
\caption{Feasible-walk comparison under the original assembly pipeline. Feasible-walk ACO is the proposed sampler for legal walks without the learned transition prior. Metrics follow the original parser for consensus quality. Higher Covered, Used, and Identity are better; lower Contigs, Breaks, Indels, and Diffs are better.}
\label{tab:feasible-aco}
\begin{tabular}{llrrrrrrrr}
\hline
Graph & Solver & Time & Covered & Used & Contigs & Breaks & Indels & Diffs & Identity \\
\hline
Minigraph & MQLib      & 5 & 47.16 & 99.98 & 4.00 & 0.00 & 0.67 & 0.00 & 98.49 \\
Minigraph & Feasible-walk ACO & 5 & 59.06 & 94.59 & 5.00 & 0.67 & 0.00 & 0.00 & 99.31 \\
\hline
\end{tabular}
\end{table}

Table~\ref{tab:feasible-aco} isolates the effect of moving feasibility from objective penalties into the sampling space. Relative to MQLib, the sampler increases coverage on minigraph-guided assemblies while trading off used sequence and continuity, which indicates that direct construction can improve the assembly objective before learned guidance is added. The continuity tradeoff motivates the learned prior as a mechanism for reshaping the feasible-walk distribution.

\subsection{Optimization quality from learned transition priors}

For the learned-prior comparison, we start from GFAs with \texttt{SC:f} depth tags produced by the minigraph pipeline. After converting the tags into copy numbers in the benchmark style, we transform the graphs into oriented optimization instances, while held-out test graphs remain separate from the online training batches. Because the base and guided ACO variants share the sampler, residual heuristic, pheromone update, objective, graph distribution, and output representation, the graph neural prior provides the isolated intervention.

Under matched objective evaluation, we test whether edge-level guidance changes the quality reached by the same search distribution. On 10 held-out query GFAs from minigraph, we evaluate MQLib, Pathfinder, feasible-walk ACO, learned-prior ACO, beam search, and A* through the same Oriented Tangle Resolution objective. Together, these comparators cover the original QUBO backend, the original biological baseline for walk search, bounded-width heuristic search, classical best-first graph search adapted to this task, and the non-neural version of the proposed sampler. Pathfinder follows its native completion rule in this wrapper, so we run it once to completion on each GFA and repeat that scored output across the budget blocks. Runtime is measured end to end for each solver wrapper, including shared construction overhead for rows that use QUBO construction. The guided solver uses the ACO setting trained on trajectories from Table~\ref{tab:dynaco-settings}.

\begin{table}[t]
\centering
\scriptsize
\caption{Measured wall-clock comparison at the optimization stage on 10 held-out minigraph test instances. The table includes the direct QUBO/MQLib baseline, Pathfinder as the biological baseline for walk search, graph-search controls adapted to the same oriented objective, and ACO variants over feasible walks. LP-ACO denotes ACO with the learned prior and FW-ACO denotes feasible-walk ACO without the learned prior. Energy, gap, and runtime are means over instances. Lower energy, lower gap, and lower runtime are better; wins count instances tied for best energy at that budget. Pathfinder is run once to completion and rescored with the same objective, so its row is repeated across budget blocks. Runtime includes the solver call plus shared QUBO construction and wrapper overhead for rows that use that construction.}
\label{tab:qubo-stage}
\begin{tabular}{lrrrrr}
\hline
Budget & Solver & Mean energy & Gap & Wins & Runtime (s) \\
\hline
5 & LP-ACO & 40.78 & 0.00 & 10 & 277.24 \\
5 & A* & 42.12 & 1.34 & 0 & 19.51 \\
5 & FW-ACO & 44.20 & 3.42 & 0 & 16.81 \\
5 & Beam search & 44.51 & 3.73 & 0 & 19.21 \\
5 & MQLib & 193.03 & 152.26 & 0 & 25.05 \\
5 & Pathfinder & 211.24 & 170.46 & 0 & 6.66 \\
\hline
10 & LP-ACO & 40.78 & 0.13 & 5 & 285.53 \\
10 & Beam search & 41.22 & 0.57 & 5 & 33.16 \\
10 & A* & 42.12 & 1.47 & 0 & 35.02 \\
10 & FW-ACO & 44.34 & 3.69 & 0 & 17.69 \\
10 & MQLib & 149.22 & 108.57 & 0 & 41.62 \\
10 & Pathfinder & 211.24 & 170.59 & 0 & 6.66 \\
\hline
180 & LP-ACO & 40.78 & 0.13 & 5 & 548.39 \\
180 & Beam search & 41.22 & 0.57 & 5 & 35.96 \\
180 & A* & 42.12 & 1.47 & 0 & 545.26 \\
180 & FW-ACO & 44.97 & 4.32 & 0 & 185.55 \\
180 & MQLib & 89.78 & 49.12 & 0 & 552.79 \\
180 & Pathfinder & 211.24 & 170.59 & 0 & 6.66 \\
\hline
\end{tabular}
\end{table}

Table~\ref{tab:qubo-stage} shows that the prior shifts the search distribution toward walks with lower copy-number loss on held-out minigraph instances. At the shortest budget, learned-prior ACO wins all 10 instances and lowers mean energy from 42.12 for A* and 44.20 for the unguided variant to 40.78. Pathfinder finishes quickly, with emitted paths that score higher under this objective. At 10 seconds and above, the guided solver and beam search split the instance wins, although the guided solver retains the lowest mean energy. Thus, the quality result is consistent across budgets, whereas the runtime result exposes a separate implementation cost: the current neural implementation is substantially slower because each instance runs the full trajectory-trained ACO loop with repeated neural-guided sampling. These results support learned transition guidance as a mechanism for solution quality, with runtime remaining an implementation tradeoff.

\subsection{Ablation: learned guidance versus control priors}

The prior ablation asks whether the learned component contributes useful graph information or merely perturbs the ACO transition distribution. To isolate this question, we evaluate three stochastic seeds on the 10 held-out minigraph optimization instances at the 5 second configured budget, using the same sampler over feasible walks and changing only the prior source and neural weight $\gamma$. A zero prior removes neural information, a shuffled prior preserves the learned multiset of logits while breaking edge identity, and a random prior injects Gaussian edge scores with the same scale as the learned logits. We report the walk-coverage proxy cost used by the feasible-walk sampler in these runs.

\begin{table}[t]
\centering
\small
\caption{Ablation of the learned prior at the 5 second configured budget on 10 held-out minigraph instances and three stochastic seeds. Rows report 30 evaluations. Gap is relative to the best condition for the same seed and instance. Lower proxy cost and gap are better.}
\label{tab:prior-ablation}
\begin{tabular}{llrrrr}
\hline
Prior & $\gamma$ & Rows & Mean proxy cost & Gap & Wins \\
\hline
Learned & 0.25 & 30 & -0.4742 & 0.0120 & 25 \\
Learned & 0.50 & 30 & -0.4742 & 0.0121 & 20 \\
Learned & 1.00 & 30 & -0.4741 & 0.0121 & 10 \\
Zero & 1.00 & 30 & -0.4828 & 0.0035 & 5 \\
Shuffled & 0.25 & 30 & -0.3755 & 0.1107 & 0 \\
Random & 0.25 & 30 & -0.3446 & 0.1417 & 0 \\
Shuffled & 1.00 & 30 & -0.0172 & 0.4691 & 0 \\
Random & 1.00 & 30 & 0.2312 & 0.7174 & 0 \\
\hline
\end{tabular}
\end{table}

In Table~\ref{tab:prior-ablation}, the learned prior is substantially better than shuffled or random score controls, so the active guidance is tied to edge identity rather than to arbitrary perturbation. The zero-prior control has the lowest mean proxy cost on this short optimization-only slice, while the learned prior at $\gamma=0.25$ wins 25 of 30 seed-instance comparisons and stays close in mean cost. This ablation therefore separates two effects: learned logits provide edge-specific guidance when a prior is active, and the paired full-assembly evaluation below tests the learned prior under the assembly metric.

To evaluate transfer beyond the held-out set with 73--79 segments, we also use a second minigraph test slice with mean segment count 100. In this scale evaluation, we use the same trained checkpoint and the same configured solver limits but change the selected GFA instances. Table~\ref{tab:qubo-scale} shows that the guided solver remains the lowest-energy method on all 12 instances across the tested limits. We interpret this as evidence for transfer across scale at the optimization stage, while preserving the runtime pattern from Table~\ref{tab:qubo-stage}: the current neural implementation has higher measured runtime than the non-neural baselines.

\begin{table}[t]
\centering
\scriptsize
\caption{Scale evaluation at the optimization stage on 12 held-out minigraph instances with mean segment count 100. All rows use the same trained checkpoint for the learned prior. Configured limit is the solver limit supplied to each method; measured runtime includes solver call and shared wrapper overhead. Lower energy and lower gap are better.}
\label{tab:qubo-scale}
\begin{tabular}{lrrccrc}
\hline
Solver & Configured limit & Rows & Mean energy & Gap to best & Wins & Measured runtime (s) \\
\hline
LP-ACO & 1 & 12 & $55.95$ & $0.00$ & 12 & $36.60$ \\
FW-ACO & 1 & 12 & $67.00$ & $11.04$ & 0 & $11.69$ \\
Beam search & 1 & 12 & $79.14$ & $23.19$ & 0 & $11.66$ \\
MQLib & 1 & 12 & $340.44$ & $284.49$ & 0 & $19.95$ \\
\hline
LP-ACO & 5 & 12 & $55.95$ & $0.00$ & 12 & $31.99$ \\
FW-ACO & 5 & 12 & $65.79$ & $9.84$ & 0 & $13.60$ \\
Beam search & 5 & 12 & $69.64$ & $13.69$ & 0 & $13.42$ \\
MQLib & 5 & 12 & $419.31$ & $363.35$ & 0 & $21.81$ \\
\hline
LP-ACO & 10 & 12 & $55.95$ & $0.00$ & 12 & $29.69$ \\
FW-ACO & 10 & 12 & $67.10$ & $11.15$ & 0 & $18.00$ \\
Beam search & 10 & 12 & $63.66$ & $7.70$ & 0 & $17.94$ \\
MQLib & 10 & 12 & $557.01$ & $501.05$ & 0 & $26.12$ \\
\hline
\end{tabular}
\end{table}

\subsection{End-to-end assembly results}

To test whether intermediate optimization gains translate to actual assembly quality, we evaluate minigraph assemblies with the original consensus parser in a paired comparison. For each seed, the same generated data-only directory is reused by every solver before path finding and consensus evaluation. We use seeds 00001--00008, with five held-out sequences per seed, for 40 assemblies per complete row, and the configured limits for path finding define the assembly conditions. The model is trained as a minigraph-specific coverage prior, fixed as a learned checkpoint, and then run through the same paired MG assembly benchmark as FW-ACO and Pathfinder. Because the paired provenance includes the native Pathfinder row, the coverage comparison uses the same data-generation runs and the same consensus parser.

\begin{table}[t]
\centering
\small
\caption{Full minigraph assembly comparison through the original consensus parser. LP-ACO denotes ACO with the learned prior and FW-ACO denotes feasible-walk ACO without the learned prior. Configured limit is the time limit for each path-finding call, not measured end-to-end runtime. Higher Covered, Used, and Identity are better; lower Contigs and Breaks are better.}
\label{tab:partial-full-assembly}
\begin{tabular}{llrrrrrrr}
\hline
Graph & Solver & Configured limit & Seqs & Covered & Used & Contigs & Breaks & Identity \\
\hline
Minigraph & Pathfinder & -- & 40 & 77.21 & 87.15 & 5.62 & 2.48 & 99.35 \\
Minigraph & FW-ACO & 30 & 40 & 77.65 & 71.01 & 1.00 & 7.60 & 98.19 \\
Minigraph & LP-ACO & 30 & 40 & 80.84 & 79.11 & 1.00 & 4.95 & 98.40 \\
\hline
\end{tabular}
\end{table}

Table~\ref{tab:partial-full-assembly} shows that the learned prior improves the full minigraph assembly comparison when it is trained and evaluated as a minigraph-specific coverage prior. At the 30 second configured limit, LP-ACO improves coverage from 77.65 to 80.84, increases used sequence from 71.01 to 79.11, and reduces breaks from 7.60 to 4.95 relative to FW-ACO on the same 40 paired assemblies. Against Pathfinder on the same paired data-only directories, LP-ACO improves the coverage target from 77.21 to 80.84 and has higher coverage on 29 of 40 assemblies. Pathfinder remains stronger on used sequence, break count, and identity, so the current evidence positions LP-ACO as a coverage-specialized learned prior rather than a uniform replacement for all consensus metrics. These results show that learned transition guidance can transfer from trajectory-level optimization to end-to-end consensus quality when the prior is trained for the same minigraph coverage objective used by the assembly benchmark.

\begin{table}[t]
\centering
\small
\caption{Coverage-aligned learned-prior ablation on the same 40 paired minigraph assemblies. Both rows use the neural ACO path finder and the same trained checkpoint; the zero-prior row loads the checkpoint but sets the edge prior to zero during sampling. Higher Covered, Used, and Identity are better; lower Contigs and Breaks are better.}
\label{tab:coverage-prior-ablation}
\begin{tabular}{llrrrrrrr}
\hline
Graph & Solver & Configured limit & Seqs & Covered & Used & Contigs & Breaks & Identity \\
\hline
Minigraph & LP-ACO zero prior & 30 & 40 & 77.84 & 70.26 & 1.00 & 7.70 & 98.18 \\
Minigraph & LP-ACO & 30 & 40 & 80.84 & 79.11 & 1.00 & 4.95 & 98.40 \\
\hline
\end{tabular}
\end{table}

Table~\ref{tab:coverage-prior-ablation} aligns the prior ablation with the end-to-end coverage metric rather than only with the optimization proxy. With the same paired assemblies and the same neural ACO implementation, enabling learned edge logits improves mean coverage by 3.00 points over the zero-prior control, improves used sequence by 8.85 points, and reduces breaks by 2.75. At the individual-assembly level, the learned prior has higher coverage in 19 cases, ties in 9, and has lower coverage in 12. This confirms that the active prior contributes to the reported coverage gain, while the mixed per-instance outcomes motivate paired means and win counts rather than a uniform-dominance claim.

\section{Limitations}

Our optimization evidence covers held-out minigraph instances in the 73--79 segment range, with an additional scale evaluation at 100 segments, while the end-to-end evidence covers the reported minigraph assembly comparison on 40 sequences. Broader scale claims require repeated timing measurements across additional graph sizes and an evaluation path that separates QUBO construction overhead from solver time. The current neural implementation can improve held-out optimization scores but is slower in measured wall-clock time than beam search and the base ACO sampler.

Our method inherits upstream uncertainty from graph construction, read-to-graph annotation, and copy-number estimation. Because the sampler chooses among paths present in the pangenome graph, sample-specific sequence represented outside that graph is beyond the current action space, and the residual heuristic may over-favor long high-copy nodes when estimates are biased. The reported guided experiments focus on the minigraph annotation route from the broader benchmark; GraphAligner, vg Giraffe, and kmer2node provide natural transfer tests. Although we train the prior on minigraph assembly instances, the dense-QUBO evaluation path restricts training to a feasible subset of that distribution; removing this subset restriction requires sparse QUBO construction, while a direct path for training on walk energy would avoid materializing the full matrix indexed by time. Finally, energy and assembly metrics remain distinct evaluation layers, making stochastic stability an important target for future work.

\section{Conclusion}

We solve Oriented Tangle Resolution in the space of feasible oriented walks rather than in a binary assignment space constrained by penalties, making construction of legal walks the search primitive. The ACO instantiation provides residual-copy heuristics, termination handling, and pheromone feedback for this pangenome-specific action space while keeping the objective and consensus parser of the original pipeline and avoiding the invalid time-indexed assignments used by QUBO/MQLib. We then train a neural prior on trajectories to supply graph-conditioned guidance over legal successor choices. This prior provides edge-specific guidance relative to shuffled and random controls, although the zero-prior control remains competitive on the short optimization proxy. In full minigraph assembly, learned guidance improves coverage, used sequence, and structural continuity over the non-neural sampler on paired held-out assemblies. The result is an optimization interface over feasible walks for pangenome path finding, with neural priors as a data-driven extension whose runtime, transfer across annotation sources, integration of edge evidence, and stochastic stability define the next implementation targets.

\bibliographystyle{unsrtnat}
\bibliography{references}

\end{document}
